\begin{figure}[h]
\centering
\begin{tikzpicture}[scale=.6, thick]
    % First plot: 5 linear regions
    \begin{scope}[shift={(0, 0)}]
        \draw[->] (-0.5, 0) -- (7, 0) node[right] {};
        \draw[->] (0, -1.5) -- (0, 1.5) node[above] {};
        \draw[wp-text, thick, domain=0:6.28, samples=100, smooth] plot (\x, {sin(\x r)});
        \draw[wp-primary, thick] (0, 0) -- (1.57, 1);
        \draw[wp-primary, thick] (1.57, 1) -- (3.14, 0);
        \draw[wp-primary, thick] (3.14, 0) -- (4.71, -1);
        \draw[wp-primary, thick] (4.71, -1) -- (6.28, 0);
        \node[above] at (6.28, 0) {\textcolor{wp-text}{$\sin(x)$}};
        \node[above] at (3.14, 1) {\textcolor{wp-text}{5 linear regions}};
        \foreach \x/\y in {0/0, 1.57/1, 3.14/0, 4.71/-1, 6.28/0} \fill[wp-primary] (\x, \y) circle (1pt);
    \end{scope}
    % Second plot: 20 linear regions
    \begin{scope}[shift={(9, 0)}]
        \draw[->] (-0.5, 0) -- (7, 0) node[right] {};
        \draw[->] (0, -1.5) -- (0, 1.5) node[above] {};
        \draw[wp-text, thick, domain=0:6.28, samples=100, smooth] plot (\x, {sin(\x r)});
        \foreach \i in {0,...,19} {
            \pgfmathsetmacro{\xstart}{\i * 6.28 / 20}
            \pgfmathsetmacro{\xend}{(\i + 1) * 6.28 / 20}
            \pgfmathsetmacro{\ystart}{sin(\xstart r)}
            \pgfmathsetmacro{\yend}{sin(\xend r)}
            \draw[wp-primary, thick] (\xstart, \ystart) -- (\xend, \yend);
        }
        \node[above] at (6.28, 0) {\textcolor{wp-text}{$\sin(x)$}};
        \node[above] at (3.14, 1.3) {\textcolor{wp-text}{20 linear regions}};
        \foreach \i in {0,...,20} {
            \pgfmathsetmacro{\x}{\i * 6.28 / 20}
            \pgfmathsetmacro{\y}{sin(\x r)}
            \fill[wp-primary] (\x, \y) circle (0.5pt);
        }
    \end{scope}
\end{tikzpicture}
\caption{Approximating $\sin(x)$ with piecewise linear functions: more hidden units (linear regions) give better approximation, illustrating the universal approximation theorem.}
\label{fig:UAT_fixed}
\end{figure}